% ID: SEL-D0047-OB-E02
% ESTADO: cotejada. Segunda pasada visual y borrador compilado revisados.
% Fuente: Selectividad/sel_2014_matematicas_aplicadas/Matematicas_CCSS_SEPT.pdf
% SHA-256: ee5bd8e23465d55ed58339e7d9adf3d58423ac2966b53d92e39a2f9b35558ab1
% Página: 2; opción B; ejercicio original 2.

Sea la función $f(x)=-x^2+px+q$.
\begin{enumerate}[label=\alph*)]
\item \textbf{(1.5 puntos)} Calcule los valores que deben tener $p$ y $q$ para que la gráfica de la función $f$ pase por el punto $(-4,-5)$ y presente un máximo en el punto de abscisa $x=-1$. Determine el valor de $f(x)$ en ese punto.
\item \textbf{(1 punto)} Represente la gráfica de $f$ para $p=2$ y $q=-1$ y halle la ecuación de la recta tangente a esta gráfica en el punto de abscisa $x=-2$.
\end{enumerate}
